\section{Methodology}
\label{sec:methodology}

\subsection{Overview}

Our study addresses a fundamental question in LLM code repair: which feedback signal---execution test results or pylint/mypy static analysis warnings---produces larger pass@1 improvement when both are given the same inference compute budget? The iso-compute constraint is the methodological core: without it, apparent feedback quality differences may simply reflect differences in the number of tokens spent on repair.

\subsection{Experimental Conditions}

We compare three conditions on HumanEval (164 problems) and MBPP (378 problems):

\textbf{Condition A: No-feedback baseline.} Single-pass greedy decoding. No repair. All token budget $B$ spent on initial generation.

\textbf{Condition B: Pylint/mypy iterative repair.} After initial generation, failing solutions receive pylint and mypy output formatted as structured feedback. The LLM re-generates given the original problem + failed solution + feedback. Repair continues until the token budget is exhausted.

\textbf{Condition C: Execution test feedback iterative repair.} After initial generation, failing solutions are executed against the benchmark test suite. The error output (exception type, traceback, expected vs.\ actual values) is formatted as structured feedback. The LLM re-generates given the original problem + failed solution + execution feedback.

\textbf{Key invariant:} All conditions use Llama 3.1 8B Instruct, greedy decoding (temperature=0, seed=42), and token budget $B$=1000 output tokens per problem.

\subsection{Iso-Compute Token Budget}

The token budget $B$=1000 output tokens per problem is the experimental unit of compute. Repair stops when the remaining budget falls below 50 tokens or when the problem passes. At $B$=1000, with \texttt{max\_tokens}=512 per round and initial generation consuming approximately 400--512 tokens, most problems complete one repair round before budget exhaustion. Per-round trajectory data confirms that rounds 2--3 contribute near-zero incremental improvement.

\subsection{Feedback Signal Design}

\textbf{Pylint/mypy feedback format.} We run \texttt{pylint --output-format=text} (default configuration, all categories) and \texttt{mypy} on each failing solution. Output is parsed and formatted as structured prompt context.

\textbf{Execution feedback format.} We execute the failing solution in a sandboxed subprocess (15-second timeout) against the benchmark's test assertions. The exception type, traceback ($\leq$512 chars), and first failing assertion are formatted as structured feedback.

\subsection{Statistical Analysis}

\textbf{Primary test: McNemar's test} on the paired 2$\times$2 contingency table (condition-passes $\times$ condition-fails), $\alpha$=0.05.

\textbf{Effect size: $\Delta_{\text{pass@1}}$} = pass@1(condition) $-$ pass@1(no-feedback baseline).

\textbf{Bootstrap CIs:} 95\% confidence intervals, 10,000 samples, seed=42.

\subsection{Mechanism Study: Pylint Coverage Analysis}

For the 64 HumanEval baseline failures, we run pylint and mypy and record all flags by category (E: Error, W: Warning, C: Convention, R: Refactor, I: Information). We compute total coverage, functional coverage (E+W), and the category distribution of all flags.

\subsection{Model and Infrastructure}

\begin{table}[h]
\centering
\caption{Model and infrastructure parameters.}
\label{tab:infra}
\begin{tabular}{ll}
\toprule
\textbf{Parameter} & \textbf{Value} \\
\midrule
Model & Llama 3.1 8B Instruct \\
Backend & vLLM v0.10.1.1 (bfloat16) \\
Hardware & 5$\times$ H100 NVL \\
Token budget $B$ & 1000 output tokens per problem \\
Max repair rounds & 3 (effectively 1 at $B$=1000) \\
Execution timeout & 15 seconds \\
\bottomrule
\end{tabular}
\end{table}
